\documentclass{ktuthesis}

\title{Baigiamojo projekto pavadinimas}
\author{Studentė Studentaitė}
\supervisor{lekt. Vadovas Vadovauskas}
\reviewer{doc. Recenzentas Recenzaitis}
\authorgenitive{Studentės Studentaitės}

\usepackage{lastpage}

\begin{document}
  \maketitle

  \thispagestyle{empty}
  Studentaitė, Studentė. Baigiamojo projekto pavadinimas. Bakalauro studijų baigiamasis projektas / vadovas lekt. Vadovas Vadovauskas; Kauno technologijos universitetas, Informatikos fakultetas.

  Studijų kryptis ir sritis (studijų krypčių grupė): Informatikos mokslai, Programų sistemos.

  Reikšminiai žodžiai: .........(įrašykite).

  Kaunas, \the\year. \pageref{LastPage} p.p.

  \begin{center}
    \textbf{Santrauka}
  \end{center}

  Lorem ipsum dolor sit amet, eam ex decore persequeris, sit at illud lobortis atomorum. Sed dolorem quaerendum ne, prompta instructior ne pri. Et mel partiendo suscipiantur, docendi abhorreant ea sit. Recteque imperdiet eum te.

  Eu eum decore inimicus consetetur, cu usu habeo corpora intellegam. Ut antiopam efficiendi deterruisset sit. Mel sint eirmod id, qui quot virtute id, dolor nemore forensibus usu id. Fugit dolore voluptatum cu vim. An vix veniam graecis insolens, sit posse iusto id. Ut vim ceteros percipit, id quo ubique recusabo, eum sint lucilius ea. In sumo inani numquam has.

  \clearpage

  \thispagestyle{empty}
  Studentaitė, Studentė. Title of the Final Degree Project. Bachelor's  Final Degree Project / supervisor lect. Vadovas Vadovauskas; nformatics Faculty, Kaunas University of Technology.

  Study field and area (study field group): Computer Sciences, Software Systems.

  Keywords: .........(type here).

  Kaunas, \the\year. \pageref{LastPage} p.p.

  \begin{center}
    \textbf{Summary}
  \end{center}

  Lorem ipsum dolor sit amet, eam ex decore persequeris, sit at illud lobortis atomorum. Sed dolorem quaerendum ne, prompta instructior ne pri. Et mel partiendo suscipiantur, docendi abhorreant ea sit. Recteque imperdiet eum te.

  Eu eum decore inimicus consetetur, cu usu habeo corpora intellegam. Ut antiopam efficiendi deterruisset sit. Mel sint eirmod id, qui quot virtute id, dolor nemore forensibus usu id. Fugit dolore voluptatum cu vim. An vix veniam graecis insolens, sit posse iusto id. Ut vim ceteros percipit, id quo ubique recusabo, eum sint lucilius ea. In sumo inani numquam has.
  \clearpage

  \tableofcontents
  \clearpage

  \listoftables
  \clearpage

  \listoffigures
  \clearpage

  \begin{center}
    \textbf{Santumpų ir terminų žodynas}
  \end{center}

  \textbf{Santrumpos:}

  Doc. – docentas;

  Lekt. – lektorius;

  Prof. – profesorius.

  \textbf{Terminai}:

  Saityno analitika – lorem ipsum dolor sit amet, eam ex decore persequeris, sit at illud lobortis atomorum. Sed dolorem quaerendum ne, prompta instructior ne pri. Et mel partiendo suscipiantur, docendi abhorreant ea sit. Recteque imperdiet eum te.

  Tinklaraštis – lorem ipsum dolor sit amet, eam ex decore persequeris, sit at illud lobortis atomorum. Sed dolorem quaerendum ne, prompta instructior ne pri. Et mel partiendo suscipiantur, docendi abhorreant ea sit. Recteque imperdiet eum te.

  Beje, darbe rekomenduojame pateikti tik svarbesnes ir mažiau žinomas santrumpas bei terminus (tarkime tokių santrumpų kaip HTML, PC, IT paaiškinti nereikia)
  \clearpage

  \begin{center}
    \textbf{Įvadas}
  \end{center}

  Supažindinama su darbo specifika, aktualumu, išdėstomi tikslai bei uždaviniai, aptariama dokumento struktūra. Šiame skyriuje apie darbą kalbama abstrakčiai, nederėtų pateikti nuorodų į kitus šaltinius (1 – 2 lapai).

  \textit{Darbo problematika ir aktualumas}

  Apibrėžiama darbo problematika ir aptariamas aktualumas. Šiame poskyryje taip pat nurodoma su darbu susijusi sritis, praktinė darbo reikšmė.

  \textit{Darbo tikslas ir uždaviniai}

  Suformuluojamas pagrindinis darbo tikslas, kuris išskaidomas į kelis uždavinius (3 – 6 uždaviniai), skirtus tikslui pasiekti. Išvados dokumento pabaigoje formuluojamos uždavinių pagrindu.
    1. Pirmasis uždavinys
    2. Antrasis uždavinys
    3. Trečiasis uždavinys

  \textit{Darbo struktūra}

  Aptariama dokumento struktūra. Nurodoma kiek ir kokių skyrių dokumente yra ir kokia informacija juose pateikiama.

  \textit{Sistemos apimtis}

  Nors tai nėra tiesiogiai įvadinis sistemos aprašas, rekomenduojame būtent čia nurodyti jau realizuotos sistemos apimtį. Matai gali būti įvairūs: darbo valandos, kodo eilučių skaičius, komponentų, klasių, modulių kiekis, teikiamų paslaugų skaičius ir pan.

  \clearpage

  \section{Analizė}

  Su darbo problematika susijusios informacijos analizė (4 – 8 lapai). Skyriaus pavadinimas ir struktūra priklauso nuo baigiamojo darbo specializacijos ir pačios temos specifikos.

  \subsection{Techninis pasiūlymas}

  \subsubsection{Sistemos apibrėžimas}




\end{document}
